\documentclass[margin=10pt]{standalone}
\usepackage{tikz}
\usetikzlibrary{shapes.geometric, arrows}

\begin{document}

\tikzstyle{block} = [draw, fill=blue!10, rectangle, minimum height=3em, minimum width=4em]
\tikzstyle{sum} = [draw, fill=blue!10, circle, node distance=1cm]
\tikzstyle{input} = [coordinate]
\tikzstyle{output} = [coordinate]

\begin{tikzpicture}[auto, node distance=2cm,>=latex']
    % Nodes
    \node [input, name=input] {};
    \node [sum, right of=input] (sum) {};
    \node [block, right of=sum] (controller) {$R(s)$};
    \node [block, right of=controller, node distance=3cm] (system) {$H(s)$};
    \node [output, right of=system, node distance=2cm] (output) {};

    % Connections
    \draw [draw,->] (input) -- node {$y_r(t)$} (sum);
    \draw [->] (sum) -- node {$e(t)$} (controller);
    \draw [->] (controller) -- node {$u(t)$} (system);
    \draw [->] (system) -- node [name=y] {$y(t)$}(output);
    
    % Unity Feedback Path
    % We go from the output branch point (y) down and then back to the sum
    \draw [->] (y) -- ++(0,-2.0) -| node[pos=0.99] {$-$} (sum);
    
\end{tikzpicture}

\end{document}
